\subsection{A PSPACE decision procedure for $\GWL$}\label{app:pspace}
%

We sketch a non-deterministic PSPACE procedure to  check validity in  $\GWL$
Let $\varphi$ be a formula. As discussed in Sec.~\ref{sec_calculus}, if $\varphi\not\in\GWL$,
there exists a tree-like countermodel $\M=\stru{W,R,e}$ for $\varphi$ such that
the number of worlds is $O(\size{\varphi}^{\size{\varphi}})$.
We can also bound the number $N$ of values  needed to define $R$ and $e$:
in the worst case, the values  $e(w,p)$ and $R(w,w')$ 
must all be pairwise distinct,
for every  world $w$, $w'$ and every propositional variable $p$ occurring in $\varphi$;
accordingly,  $N$ is in  $O(\size{\varphi}^{\varphi})$.
Let $\Qcal=\{0\}\cup\{~1/k~|~1\leq k\leq N\}$;
note that the elements of $\Qcal$ can be non-deterministically generated in  PSPACE.
By the above discussion, if $\varphi\not\in\GWL$  there exists a countermodel $\M$  such that
the images of  $R$ and $e$ are contained in $\Qcal$.
To verify that such  a  countermodel can be constructed, we introduce a map $\Phi$ 
that specifies the intended values of the subformulas of $\varphi$  at each world.
More precisely,
 for every world $w$,  $\Phi(w)$ is a set of constraints  of the form $w:\a =r$, where
 $\a$ is a subformula of $\varphi$ and $r\in \Qcal$,
declaring that $e(w,\a) = r$.
The procedure starts by  setting $\Phi(w_0)=\{w_0:\varphi= r_0\}$, where $w_0$ will become the root of $\M$ and
$r_0$ is any value  non-deterministically chosen from $\Qcal\setminus\{1\}$.
 We then apply
propositional saturation steps to $\Phi(w_0)$ as long as possible. In
each step, a non-atomic constraint $w_0:\psi = r$ in $\Phi(w_0)$ is
removed and replaced by the constraints required to justify the
intended evaluation $e(w_0,\psi) = r$.
For instance, let  $\psi=\a\to\b$.
If $r=1$, we must have  $e(w_0,\a) \leq e(w_0,\b)$;
we non-deterministically pick two values $r_1$ and $r_2$ from $\Qcal$
such that $r_1\leq r_2$, and
we replace $w_0:\a\to \b=r$
with  $w_0:\a = r_1$ and  $w_0:\b = r_2$. 
Similarly, if  $r<1$  we select a value $r_1$ from $\Qcal$ such
that $r_1> r$, and we replace  $w_0:\a\to \b=r$ with $w_0:\a = r_1$
and  $w_0:\b = r$.
At the end of this saturation phase, the constraints
in  $\Phi(w_0)$ are either atomic or modal.
If $\Phi(w_0)$ contains two   constraints $w_0: \a=r_1$ and  $w_0:\a =r_2$  such that $r_1\neq r_2$,
the construction  fails. 
If $\Phi(w_0)$  contains only atomic constraints, the construction succeeds.
%hence  $\varphi\not\in\GWL$.
Otherwise, for every modal constraint  $w_0:\psi=r$ in   $\Phi(w_0)$ 
we have to define an $R$-successor of $w_0$ and continue the construction. 
More in detail,
assume that the current  modal constraint  is $w_0:\Box\a=r$ (the case  $w_0:\Diam\a=r$ is similar);
we  generate a new world $w_1$  and we  select  a value $r_1$ for $R(w_0,w_1)$ and a value
$r_2$ for $e(w_1,\a)$ such that $r_1\to r_2 = r$; thus, $w_1$ is the world  witnessing that  $e(w_0,\Box \a) = r$.
We set $\Phi(w_1) =\{ w_1:\a= r_2\}$.
We then add to  $\Phi(w_1)$ the constraints required to justify the modal constraints in $\Phi(w_0)$.
More precisely:
\begin{itemize}
\item for every $w_0:\Box \b=r'$ in  $\Phi(w_0)$, with $\b\neq \a$,  a value $r''$
such that $r_1 \to r'' \geq r'$
is  selected
  and  $w_1:\b=r''$ is added to  $\Phi(w_1)$;


\item  for every $w_0:\Diam \b=r'$ in  $\Phi(w_0)$,
 a value $r''$
such that $r_1 \land r'' \leq r'$
is  selected
  and  $w_1:\b=r''$ is added to  $\Phi(w_1)$.
\end{itemize}
Note that $\Phi(w_1)$ contains fewer modal operators than the input
formula $\varphi$; accordingly, we can recursively apply the
nondeterministic procedure to check whether a model having root $w_1$
satisfying all the constraints in $\Phi(w_1)$ can be constructed.  If
the construction succeeds, we say that the constraint $w_0:\Box\a=r$ is
checked.
Note that, at each iteration, the information needed to check a constraint
can be discarded, so the space usage remains polynomial in the size of $\varphi$.
If all the modal constraints in $\Phi(w_0)$ are checked, we conclude that
a countermodel for $\varphi$ can  be built, hence $\varphi\not\in\GWL$.
If all the non-deterministic attempts to build a countermodel fail,
we conclude $\varphi\in\GWL$.


%%% Local Variables: 
%%% mode: latex
%%% TeX-master: "goedelModalLogicWitnessNonCrisp"
%%% End: 
